% ----------------------
\subsection{Teleological explanations and the teleological fallacy}
% ----------------------

	% -----
	\subsubsection{Introduction}
	% -----

		In this essay, I'll be discussing a certain way of explaining things called \textit{teleological explanation}. I'll also explore a frequently fallacious way of thinking about the gospel which is based on teleological explanations. I call this way of thinking the \textit{teleological fallacy}. This fallacy is so common and so potentially destructive that I think it needs to be treated early on if we're trying to be systematic about the gospel. It is common because it exploits a very normal heuristic for reasoning; it is destructive because it applies that heuristic in domains where we have little or no reason to believe it will be successful.

		Before getting started, I should also say that I am not the first to talk about this fallacy. Others may have spoken about it in the context of thinking about the LDS church---though I'm not aware of anyone who has---but arguments over teleology go back thousands of years, as I'll discuss below. Evolutionary biology in particular has fought some pretty bloody battles over teleology and its use in explaining biological functions.%
			\footnote{Chapters 2 and 3 of Barrow and Tipler's \textit{The Anthropic Cosmological Principle} discusses design arguments in detail. Several passages of Dennett's \textit{Darwin's Dangerous Idea} also deal with teleological explanations and reasons.}
		So I am by no means the first to raise this concern, though I rarely see it discussed among LDS believers in these terms.
	
	% -----
	\subsubsection{Design, purpose, and teleology}
	% -----
	
		If you've ever gone house-shopping or apartment-hunting, then you know what it's like to walk through dozens of places looking for one that you could be happy with. You see small and big, dirty and clean, realistic and pipe-dream. Some of those places are already built, but you've probably also toured others that are still under construction or need to be finished.
		
		When you enter an unfinished house, you're walking through the skeleton: the frame is exposed, the studs are visible along the walls, the ceiling is open, and the foundation might be showing in the basement. Even though many of the things that make a building a home aren't there yet, you're still able to recognize the building \textit{as} a home, because you understand that it was \textit{designed} to be a home. When you walk into the kitchen area, you see the empty spots between counters as spaces where the fridge and oven will go. You recognize the hole in the wall in the basement as the dryer vent; the many little holes in the ceilings as meant for electrical cables or plumbing; the planed area out back as the future patio.
		
		You understand all these features of the structure because you know that they are there for a \textit{purpose}. If you didn't know that the building was meant to be a house---or worse, if you didn't even know that the wood and concrete were meant to be a \textit{building}---than most of what you saw when walking through it would be a mystery. ``Who put all these holes in these perfectly good boards?'', you might wonder. ``Who filled up this fine pit in the ground with solid rock?'' The objects and structures which make up the building are comprehensible because you know that they were designed to perform a certain function; they were made with a certain purpose in mind.
		
		The study of such designs and purposes is called \textit{teleology}. In ancient Greek, the word \textit{telos} (\gr{τέλος}) means ``purpose, end, goal,''%
			\footnote{This word can mean lots of other things too, but the sense of ``goal'' is prominent in ancient philosophical texts.}
		so the English word ``teleology'' refers to the study of purposes. When we reason teleologically, we reason on the basis of purposes and intentions, and then we use what we know about those intentions in order to make sense of something. Teleological reasoning is what allows us to understand that the concrete, wood, and other components of an unfinished house all have a purpose. We recognize that, if there's an empty space in the kitchen, that space must be intended for something---it must have been designed that way, or else it wouldn't be there. After all, who would put a random empty space in the middle of a perfectly good kitchen? No one, since there would be no \textit{reason} for it---and that's teleology.
		
		Even when we're confronted with an object or situation that stumps us, teleological reasoning is still our fallback. That is, even when we don't know the purpose of something, we still assume that there \textit{is} a purpose, and if we knew the purpose, we would understand the object. Our commonsense knowledge of people and homes more or less compels us to believe that the builders of the house put things together in certain ways because they had certain purposes and ends in mind.
		
		That last point is key to remember. In many cases, we don't need to know the \textit{actual purposes} someone had in mind when they made something. All we really need to assume is that there \textit{were} purposes, and that's enough to get us started on a teleological path. We can start to draw conclusions on that bare assumption: as long as there were certain purposes or intentions, this ends-based reasoning can get off the ground.
		
		Reasoning teleologically is so common in so many different circumstances that you are an absolute professional at doing it, even if no one has ever pointed out exactly what you're doing. Most of the time we don't even recognize that we're reasoning teleologically at all because, in many domains of our experience, we can't even imagine \textit{not} reasoning this way. And in many domains, teleological reasoning is a spectacular success.
		
	% -----
	\subsubsection{Human behavior as a domain of teleological success}
	% -----
		
		I've already hinted at the study of human behavior as a domain where teleological reasoning is incredibly successful. Let's take a closer look at this domain, and why teleological reasoning is so well suited for it.
		
		Suppose you're taking a walk one day and you pass a man digging a hole in his front yard. His waist is already below ground level, so the hole is pretty deep. You wonder what he's up to, and right away your teleological reasoning spins into gear: you start to think of all the reasons why a person would want a hole in their front yard. Maybe he's installing sprinklers? Maybe he's trying to get rid of a deeply-rooted plant? 
		
		So you ask the guy. ``Hey, what are you digging this hole for? Sprinklers?'' ``No,'' he responds, and keeps digging. You try again. ``Getting rid of a plant?'' ``No,'' he says again. So your mind starts spinning once more, this time proffering reasons that are a little stranger, but would still make sense: digging for gold, exercising, trying to get to China. These reasons are stranger because, in order to make sense, they require the man to believe some dubious things: that there is gold in his yard, that digging a big hole is a convenient way to exercise, that he really could get to China by digging. 
		
		Finally you just ask him: ``Why are you digging?'' ``I'm looking for a time capsule my father buried when I was child,'' he replies. Suddenly it all makes sense. He's been looking for something important to him, and sometimes people will go to great lengths to find things that are important to them, so he had a purpose all along---he had a reason for acting the way he did.
		
		Almost all human behavior works in the same way: motivated by \textit{reasons}. And this is the reason (no pun intended) why teleological reasoning (again, just a coincidence) works so well---because human activity is chock-full of reasons, intentions, purposes, and motivations. We have these reasons and intentions in mind when we act, and when others understand them, they can use them to make sense of what we believe about the world, what we want, and what our overarching goals might be. And when people can do that, looking at each other teleologically gives us a very powerful perspective for predicting and explaining behavior.%
			\footnote{This way of explaining human behavior is sometimes called ``psychological explanation,'' but psychological explanation also refers to the attribution of mental states more generally, not just intentions or purposes.}
			
		It's also important to realize that the teleological picture of behavior can succeed even when people are not consciously entertaining their reasons for acting. The man digging the hole in his yard might be thinking of the time capsule the entire time, reflecting on his father and childhood experiences. But since so many of our activities are routine, we don't always have our reasons for acting at the forefront of our minds---when I drive to work in the morning, I don't have to keep reminding myself, ``I'm doing this because I need the money'' (or maybe I do, depending on the day). The reasons are still there driving the actions, even if it is happening beneath the surface.
		
		This last point about hidden reasons is crucial because it allows us to extend the domain of teleological reasoning beyond human behavior and into the realm of living things more generally. Dogs, birds, fish, and other organisms act for reasons---the dog scratches at the door \textit{because} she wants to chase a squirrel outside; the bird builds a nest \textit{because} it needs a place to lay its eggs; the fish swims away from disturbances in the water \textit{because} it wants to avoid possible predators. These animals are probably not able to self-consciously reflect on their reasons for acting---certainly they cannot do so through language. But the reasons are still there, making it possible to take a teleological perspective on their behavior as well, which means we can predict and explain what they do in similar ways.%
			\footnote{Dennett calls this perspective the ``intentional stance.''}
		
		The general principle we want to draw out is the following: teleological reasoning tends to succeed when the domain we're using it in is populated by \textit{reasons}, where the ``reasons'' we're talking about are the products of a design process (more on the design process in the next section). The reasons, as products of a process, tell us why that process was set up in a certain way---that is, they tell us what the process was intended to achieve. Human behavior, and the behavior of other organisms, are domains populated and indeed overrun by reasons; therefore, teleological reasoning tends to succeed in them.
		
		Our general principle may look like a tautology, but it isn't.%
			\footnote{A ``tautology'' is a statement that is true under any interpretation. For example if ``A'' represents any statement whatever, then ``A or not-A'' is a tautology, because that whole statement is \textit{always} true---either statement A by itself is true, or it isn't true, and so it will always be the case that A is either true or false. In general, tautologies are not very interesting exactly because they are always true, and so they don't seem to communicate much to us; they're a bit empty, or just trivially true.}
		Whether a certain domain is populated by reasons is a discovery we have to make about that domain, involving whether there was some process of design. If and when we are able to make that discovery, then we have good evidence to believe that teleological reasoning will succeed there. But sometimes we make the wrong assumption about the presence of reasons, and this leads to failures of teleology, as we'll see shortly.
		
		Before the failures, though, let's look at a case of reasons involving design, where we have to actually choose between two different sets of reasons. Each set has very different consequences for both how we apply our teleological reasoning, and what that reasoning means more broadly.
		
	% -----
	\subsubsection{Organism diversity: a choice between two designs}
	% -----
	
		One of the most striking things about the natural world is how well-suited different organisms are to the life they live. The beaks of birds are exactly fitted to the things they eat; the coats of cold-weather mammals are thick and warm; sharks are streamlined to glide through the water with little resistance in order to catch prey efficiently.
		
		And I've only given a few examples---it's even \textit{more} striking that these more or less optimal features are present in \textit{every} organism, no matter where they live or how they behave. And since there are millions of different species and great variety in sub-species as well, some sort of optimization process must have occurred for each and every one of them, making them into the well-built eating and reproducing machines that they are.
		
		Optimization leads to reasons that we can grasp and use in our teleological picture of the world. Why is a polar bear's coat so thick and warm? The reason is that he lives where it is very cold, and needs to be protected from the weather. The reason a shark is so streamlined is that she must catch prey that in most cases is smaller and quicker than her, and so she has to have a fighting chance at getting her fill. By assuming that there has been some sort of optimization for an organism's features---even if that process wasn't completely optimal---we enter a domain populated by reasons and can therefore rely on teleological logic to guide us.
		
		What's interesting about biology is that there are two very clear and very different ways of explaining why these reasons exist. The first option is to believe that God designed each species, or even each organism, and that God's providential design explains why different organisms are so well-suited to their environment.
		
		A number of different things make the God-option for biology very compelling. One is that we often explain other instances of \textit{design} in other domains by appealing to a process of intelligent or mind-guided design. That's how we explain the intricate design of my house, for example, or of my computer or my car. Those objects were designed by minds with intentions, so we explain their features on that basis. So why not explain biological phenomena in the same way? God as intelligent designer could have optimized the structure and behavior of each organism to match the sort of world they would live in.
		
		The second option, of course, is Darwin's theory of natural selection, or at least some version of it. On this view, we explain the apparent design of the biological world by a process of descent with modification: mutations cause organisms to change their structure or behavior, and when these changes are beneficial they get retained, and over time the changes add up to make it appear as though there was a mind designing the polar bear, when in reality all that happened was that any polar bear-like thing without a warm coat just didn't reproduce enough for his genes to carry on.
		
		From both of these options it follows that there are reasons why organisms are how they are, which means that teleological reasoning is going to work either way. One of Darwin's great achievements was to show how there could be reasons \textit{without a mindful designer}---and so how teleological reasoning could succeed even when there was no \textit{intelligence} or \textit{mind} guiding the reasons-producing process you want to think teleologically about.%
			\footnote{You might call it \textit{Darwin's Dangerous Idea}.}
		
		Thus we're faced with a choice between two options that both give explanations and both provide reasons. How do we decide between them? I'll have a bit more to say about this in a later essay, but the answer is that we have to choose the option which fits better with other things we think we know. In this case, the ``other things'' include stuff we've learned from geology, molecular biology, and other fields. Both of the options can be made to fit these other theories and facts, but one appears to fit them more naturally and with less stretching.
		
		Biology, then, is a domain so fully populated with reasons that we actually have to decide between two competing accounts of where the reasons come from. Either way, the reasons are the product of a design process, and their presence secures our use of teleological reasoning about both the structure and behavior of living things.
		
	% -----
	\subsubsection{Teleological failures}
	% -----
	
		Other domains of the natural world are not as lucky as biology, and don't enjoy such a wealth of reasons.
		
		Let's take geology as an example---in particular, the attempt to explain why a certain geological feature, like a mountain range, is the way it is. As far as we can tell, there doesn't seem to be any \textit{design}, either mindful or mindless, involved in the creation and shaping of a mountain range. Instead, the explanation has to do with forces inside the earth, plates and plate tectonics, and so on. None of those processes and events are directed toward any particular purpose, and they aren't being controlled by an intelligence. They just happen, and happen to create features like a mountain range along the way.
		
		On the other hand, it's really easy to see how you \textit{could} take a teleological perspective on why the earth shows certain surface features. You might think that God created this valley \textit{so that we could live in it}; thus God either formed the mountains himself or perhaps directed the geological processes in order for them to form the mountains in a certain way. If either of those things were true, then a teleological mode of explanation would be successful in geology. You could walk through the natural world as though you were walking through a half-finished house, seeing designs and purposes and reasons in the placement of different objects. The designs and purposes and reasons would then allow you to give certain explanations and reach certain conclusions teleologically.
		
		Teleological reasoning fails in the domain of geology for at least two reasons: (1) it doesn't fit well with what else we know, and (2) it leads to a bunch of unanswered questions about why God made a volcano explode, why he made such a crappy spot of land, and so on. If you involve God in one explanation, it seems like you might as well involve him in every explanation, since you don't have a principled way of restricting his activities to the small range of things you set out to explain in the first place.
		
		Another domain where teleology fails is celestial mechanics, or the study of planetary and other orbits in our solar system. Many early theories of celestial mechanics (as in a couple thousand years ago or so) made a very crucial assumption: that the earth was at the center of the universe. Some of them made this assumption on the basis of God's design, since for them it was obvious that God would have set the earth at the center of everything else---it's the most important thing in the universe, after all!
		
		Once you make the earth-at-the-center assumption, you have to work to fit everything else in around that---so the moon, the sun, Mars, Jupiter, the stars, and everything else are all going to orbit around the earth, with the earth stationary at the center. It turns out that models which assume the earth is at the center---we can call these ``Ptolemaic'' models, after Ptolemy, an ancient astronomer whose model was very influential---can be very successful at explaining the apparent motions of the sun and objects in the night sky.
		
		For a long time, it looked as though these design-based models of the solar system were going to stick around forever. But there were small differences in the predictions made by these models and the actual observations. Eventually a sun-centered model came around---the ``Copernican'' model, after Copernicus, who gave the first thorough account of one---which, over time, were seen to fit the observed data better than the Ptolemaic models.
		
		So the battle between Ptolemaic and Copernican models of the solar system was not just a choice between which object to put at the center of everything, the sun or the earth. It was also a battle over the teleological basis of celestial mechanics, and of all astronomy in general: was it a domain in which reasons, as products of a mindful design process, could be used to explain things? If yes, then it would still be legitimate to invoke God's purposes and intentions when explaining what we see. If not, astronomy would lose the anchor of design and become a domain without intelligent reasons---a domain where teleology would fail.
		
		To make a multi-century story very short, Copernicus won, teleology lost, and the sun is at the center of the solar system (though definitely not at the center of the universe). The sorts of teleological approaches that were available to Ptolemy were closed off (at least to some extent) to Kepler, Descartes, Galileo, Newton, and others who followed Copernicus, and astronomy is now like geology in generally not offering purposes and designs to ground our theories.
		
		% not sure if I need the below paragraph
		% Geology, astronomy, and many other fields of inquiry were once thought to be populated by reasons that were products of a design process. As a result, people working in these fields once offered explanations and predictions by citing these reasons. But non-teleological reasoning turned out to be more successful.
		
		The danger in using teleology within a reason-less domain is obvious: we end up \textit{imposing} design in places where there is none, and so our subsequent teleological reasoning reflects not actual reasons in the world, but instead only \textit{reasons we attribute to some domain} based on whatever \textit{apparent} design we \textit{think we notice} or even \textit{dictate}. But apparent design is not real design, and the fact that we seem to discern purpose in something is not necessarily evidence that there is purpose in it, or that it was done for certain reasons. And without purpose and reasons, teleology fails.
		
	% -----
	\subsubsection{The teleological fallacy}
	% -----
	
		At long last we come to the teleological fallacy. We've already been discussing it for a while, so no more waiting---time to make it explicit.
		
		\begin{description}
			\item[The Teleological Fallacy] The teleological fallacy is the mistake of inferring that, because something \textit{is} a certain way, it must have been \textit{designed} to be that way.
		\end{description}
		
		Like I said, we've already seen this fallacy in action. We saw it in geology: we might commit it when we try to explain some feature of the earth's surface. We first observe that a certain feature \textit{is} a certain way---say, a particular valley is full of fertile soil. Then we infer that, because the valley \textit{is} that way, it must have been \textit{designed} to be that way. And so we conclude that God must have made the valley that way, apparently for us. That's the teleological fallacy---simple, right?
		
		We also saw the fallacy in astronomy, and it crops in evolutionary biology as well, when you have to decide whether a certain feature of an organism is a result of selection, or is just a spandrel.%
			\footnote{A ``spandrel'' is a feature that arises not as an adaptation or something that is selected for, but instead as a result of other organizational changes in an organism. See Gould and Lewontin 1979, ``The spandrels of San Marco and the Panglossian paradigm: a critique of the adaptationist programme.''}
		Other fields of inquiry face similar challenges.
		
		There is one place, however, that is much more susceptible to the teleological fallacy than most, and that's the gospel. In reasoning about the gospel, we walk a narrow bridge suspended over a bed of ravenous teleofallacygators,%
			\footnote{Unfortunately, not real.}
		because we're dealing with a domain that really does have design, purposes, and reasons. But since we don't always know \textit{where} those purposes and reasons are, it's tempting to see design everywhere---permeating the gospel world with reasons---and thereby to help ourselves to teleological explanations for whatever we observe. As I will suggest, though, many such explanations are spurious, because the inferences leading to them have committed the teleological fallacy.
		
		From here on out, then, let's focus on the gospel teleological fallacy. We can state it like this:
		
			\begin{description}
				\item[The Gospel Teleological Fallacy] The gospel teleological fallacy is the mistake of inferring that, because something \textit{is} a certain way in the gospel, God must have \textit{designed} it to be that way. 
			\end{description}
			
		Obviously this version of the fallacy is similar to its non-gospel counterpart. But I think that stating the gospel version forces us to confront more directly the sorts of things that we are liable to commit this fallacy about.

		For example, here is a list of what we might call ``features'' of the gospel, or observations we could make about the way the gospel is:
		
			\begin{compactitem}
				\item The veil of forgetfulness
				\item Men bearing the priesthood, but not women
				\item Sealings between men and women
				\item Humans made in the image of God
				\item Children being considered as innocent before a certain age
				\item Temple and ordinance symbolism
					% think of the gestures in the temple, baptism being a "burial," and so on---all these are contingent on human practices
			\end{compactitem}
		
		Many people in the LDS church---most, I think---would treat these observations like a list of features in a house. Since the features are clearly present, they must have been designed that way. After all, if they weren't designed that way, why would there be there? God would have just done something different; he must have set things up in the way he did for a reason.
		
		I think that's wrong, at least when it comes to the gospel ``features'' on that list, and probably many more. For these aspects of the gospel, at least, it is not that God \textit{designed} them this way, as a deliberate choice among several options. Rather, the explanation for them is that \textit{God \textbf{had} to set things up that way, because of the physical constraints placed on him in implementing the gospel on this earth}.
		
		Let's take the veil of forgetfulness as an example to see what I mean (I will cover the veil in greater detail later on). The ``veil of forgetfulness,'' or just the ``veil'' to most people, is the supposed barrier that exists between us and (1) our prior knowledge of and familiarity with God, before coming to this earth, and (2) the disembodied spirits which currently inhabit this earth. The LDS claim is that we lived with God before coming to this earth, but the manifest fact is that we don't remember it. How do we explain the lack of memories, the lack of knowledge, of this previous life? The standard answer is that God placed a ``veil'' over our minds upon coming to this earth, so that we cannot remember our previous experiences. The veil somehow prevents us from accessing experiences that we really did have and, apparently, ought to be able to remember, were it not for the veil.
		
		The teleological interpretation of the veil then has to go one step farther and say, since God clearly has placed the veil over our minds, he must have had good \textit{reasons} for doing so. And then we make up reasons for it---usually, that mortal life is more of a test if you don't remember your experiences before you were born on earth.
		
		Beyond the fact that that reason is a very bad one for God to set up the veil (why wouldn't God want to increase our chances of success as much as possible?), the teleological explanation for the veil also ignores the reality of human embodiment---it ignores the fact that humans think, remember, and speak by using a nervous system. And creating a nervous system from the ground up, as every human mother starts to do when she conceives, is expensive and difficult. Including memories of past experiences would require that some physical influence shape the nervous system in the womb and early childhood to carry a record of each individual's pre-earth life. Such a mechanism would have to be customized for each person, and would itself have to carry the record of that person's pre-earth experiences. When you lay it out this way, and actually look at what would be involved in eliminating the ``veil,'' I think the appeal of the teleological explanation seriously diminishes.
		
		On the other hand, on a less teleological way of thinking, there is a much simpler explanation for the veil. The phenomenon of the veil and our forgetting past experiences come about because of the constraints involved in building a human body that can think by means of natural processes. There simply is no physical mechanism that could leave a record in every human nervous system of each individual's pre-earth experiences. The explanation for the ``veil,'' therefore, is just that it is the expected result of creating human bodies through biological processes; it is false that God caused us to forget our previous experiences because of some purpose the forgetting would serve. Rather, we \textit{had} to forget, and that constraint wasn't one that was chosen or superable by God.
		
		The explanation I just gave for the veil is called a ``naturalistic explanation.'' I will be giving more of them, and will write an essay specifically dedicated to them. Naturalistic explanations do not impute design where there may not be any, or where we have reason to believe there isn't. The absence of design is one of the reasons many people don't like naturalistic explanations for features of the gospel, but as I will argue later, naturalistic explanations ought to be preferred over teleological explanations, in most cases, when we can get them.
		
		In fact, one way of viewing this entire series of essays is as an attempt to solve this exact problem: how do we reconcile the apparent purpose-driven nature of the gospel, alongside related scriptural and prophetic commentary, with the rest of what we know about the world? I believe a satisfying reconciliation is possible, and will try to make the case for it.